\chapter{Parallel reduced order modeling for digital twins using high-performance computing workflows (Article 2)}
\label{chap:article2}

\section*{Preface}

This chapter reproduces the peer-reviewed journal article published in \emph{Computers and Structures}~\cite{Ares2024}, which addresses the integration of PROM techniques with HPC workflows for industrial applications. Building on the methodologies introduced in earlier chapters, this work focuses specifically on overcoming the computational bottlenecks that arise during the offline phase of PROM construction when scaling to large, complex systems.

This contribution represents the main intrusive reduced-order modeling outcome developed within the \href{https://www.eflows4hpc.eu/}{eFlows4HPC} framework, as part of a broader European research pathway that also includes the \href{https://cordis.europa.eu/project/id/800898}{ExaQUte} and \href{https://cordis.europa.eu/project/id/946009}{EdgeTwins} projects. These initiatives provided the theoretical and technological foundations for enabling scalable, real-time digital twin applications through a combination of exascale computing, projection-based ROMs, and edge deployment strategies. The results presented here adapt and extend this foundation to a full high-fidelity industrial use case, demonstrating how an intrusive PROM workflow can be embedded within an HPC-enabled pipeline.

The development of this workflow was made possible through close collaboration between CIMNE/UPC, Universitat Politècnica de València (UPV), Siemens, Barcelona Supercomputing Center (BSC), and partners such as Inria (mesh optimization) and SISSA (non-intrusive ROM methodologies).

The chapter is structured around the following main contributions:
\begin{itemize}
    \item The development of a complete HPC-enabled workflow for intrusive PROM construction, including parallelization of training simulations, distributed SVD algorithms, and the partitioned implementation of the ECM hyper-reduction across computational subdomains, facilitating efficient parallel deployment.
    \item The validation of the proposed workflow through a multiphysics case study involving motor thermal dynamics, demonstrating its scalability, performance, and practical applicability for digital twin scenarios.
    \item The demonstration of an end-to-end Workflow-as-a-Service strategy for industrial deployment, combining HPC, cloud, and edge computing environments (e.g., through Functional Mock-Up Units (FMUs)).
\end{itemize}

The overall workflow is illustrated in Figure~\ref{fig:rom_workflow}, which situates the main intrusive pipeline alongside the mesh optimization and non-intrusive developments carried out in parallel within the same framework.

Overall, this chapter shows that under real-world industrial conditions, even highly efficient projection-based methods require HPC resources for the offline training phase, particularly when handling large-scale multiparametric studies, complex physics, and high-dimensional full-order models. By leveraging HPC workflows and integrating expertise from multiple research partners, the PROMs developed here become viable tools for real-time industrial applications, bridging the gap between high-fidelity simulation and scalable digital twin deployment.

\begin{figure}[h!]
  \centering
  \includegraphics[width=0.95\textwidth]{Figures/Workflow.png}
  \caption{Overview of the ROM workflow developed within the eFlows4HPC framework, highlighting the intrusive PROM path, mesh optimization stage (Inria), and non-intrusive ROM developments (SISSA).}
  \label{fig:rom_workflow}
\end{figure}

\includepdf[
  pages=-,
  pagecommand={},
  scale=0.92,
  delta=0 0,
  offset=32pt 0pt,
  turn=false
]{Papers/Parallel_Reduced_Order_Modeling_for_Digital_Twins_using_High_Performance_Computing_Workflows.pdf}

